\problema{Task Scheduler}{621}{src/04-heaps-topk/621-task-scheduler.cpp}{M}

Nos dan un arreglo de tareas (cada una representada por una letra de la
\texttt{A} a la \texttt{Z}) y un entero \texttt{n}: el número mínimo de
unidades de tiempo que deben pasar entre dos ejecuciones de la
\emph{misma} tarea (su ``enfriamiento'', o \emph{cooldown}). En cada unidad
de tiempo se puede ejecutar una tarea o quedarse inactivo. Hay que devolver
el número mínimo de unidades de tiempo necesarias para ejecutar todas las
tareas respetando el enfriamiento.

\paragraph{La idea, con un ejemplo de la vida real.} Imagina una lavadora
que, después de lavar una prenda de cierto color, necesita esperar
\texttt{n} ciclos antes de volver a lavar otra prenda del mismo color (para
evitar que se manchen entre sí). Si tienes muchas prendas del mismo color,
conviene lavar primero las que \textbf{más} se repiten, porque esas son las
que más tiempo de espera van a necesitar acumulado; dejarlas para el final
solo alarga más el proceso total. Por eso en cada turno conviene elegir,
entre las prendas disponibles en ese momento, la del color con más prendas
pendientes: eso es justo lo que resuelve un \emph{max-heap} de conteos.

Las prendas que se acaban de lavar no están disponibles de inmediato: se
guardan aparte, en una cola de ``enfriamiento'', hasta que les toque volver
a estar disponibles.

\begin{center}
\ttfamily
\begin{tabular}{l}
tasks = [A,A,A,B,B,B], n = 2\\
horario: A, B, idle, A, B, idle, A, B -> 8 unidades\\
\end{tabular}
\end{center}

\paragraph{Cómo lo resuelve el código, paso a paso.}
\begin{enumerate}
  \item Contamos cuántas veces aparece cada tarea con un mapa de
        frecuencias, y metemos esos conteos (solo el número, no la letra)
        en un max-heap.
  \item Usamos una cola de enfriamiento donde guardamos pares (conteo
        restante, momento en que puede volver a estar disponible).
  \item En cada unidad de tiempo: primero revisamos si el frente de la cola
        de enfriamiento ya puede regresar al heap (si su momento de
        regreso coincide con el tiempo actual), y si es así, lo regresamos.
  \item Luego, si el heap no está vacío, sacamos el conteo más alto, le
        restamos uno (la ejecutamos), y si todavía le queda conteo
        pendiente, la mandamos a la cola de enfriamiento con su nuevo
        momento de regreso (\texttt{tiempo + n + 1}).
  \item Si el heap estaba vacío en ese instante pero aún quedan tareas en
        enfriamiento, esa unidad de tiempo se pierde como inactividad.
  \item Repetimos hasta que tanto el heap como la cola de enfriamiento estén
        vacíos.
\end{enumerate}

\lstinputlisting{src/04-heaps-topk/621-task-scheduler.cpp}

\complejidad{$O(T)$}{$O(1)$}

\paragraph{¿Qué significa esa complejidad?} $T$ es el número total de
tareas. Aunque cada unidad de tiempo puede requerir una operación de heap,
el heap nunca tiene más de 26 elementos distintos (uno por letra del
alfabeto), así que su costo logarítmico es una constante que no depende de
$T$; por eso el tiempo total es lineal en $T$. El espacio también es
constante, ya que como mucho guardamos 26 conteos distintos sin importar
cuántas tareas haya en total.

\paragraph{Consejos para la entrevista (Amazon).} Este problema es un
ejemplo clásico de estrategia \emph{greedy} (voraz): en cada paso se toma la
decisión que parece mejor localmente (ejecutar la tarea con más conteo
restante), y hay que poder justificar por qué esa elección local nunca
perjudica el resultado global. También existe una fórmula matemática
cerrada basada en la frecuencia máxima y cuántas tareas empatan en esa
frecuencia máxima, que da la respuesta sin simular nada; mencionarla es un
plus, pero la simulación con heap y cola de enfriamiento es más fácil de
razonar y de explicar en la pizarra, y generaliza mejor si el entrevistador
cambia las reglas del problema a mitad de la conversación.
